\clearpage

\section {Summary of changes compared with the Run 1+2+3 analysis in T2K-TN-141(v8) }
\label{sec:changes}

The following changes have been implemented compared with the Run 1+2+3 $\nu_{\mu}$-disappearance
analysis described in \cite{OurRun123DisappAnalysis3Flavour} :

\subsection{Changes in construction of the nominal single-ring $\mu$-like reconstructed energy spectrum p.d.f}

\begin{itemize}
\item The ND280 $\nu_{\mu}$ group updated the division of reconstructed ND280 CC events. For the 2013 analyses, these events 
are divided into CC0$\pi$, CC1$\pi$ and CC other samples (whereas for the Run 1+2+3 analysis they were divided into CCQE and 
CCnonQE samples). Pion tracks are identified using ionisation energy losses in the TPC or FGD, or a pion can be identified 
by the presence of a Michel electron in FGD1. The CC0$\pi$ and CC1$\pi$ samples have exactly 0 and 1 track consistent with 
a $\pi^{+}$ respectively. The CC other sample consists of events with more than 1 track consistent with a $\pi^{+}$ or at least 
one negative track or with a track consistent with an electron \cite{ND280Inputs}. These three samples were used as an input to 
the BANFF fit of the ND280 data to find the best-fit correlated flux
and cross-section parameters \cite{BANFF2013}; these parameters (taken from \verb+postfit_banff_2013_disappearance_v4.root+) were used in 
turn to calculate the nominal predicted single-ring $\mu$-like reconstructed energy spectrum in SuperK.
\item An updated SuperK MC was used in this analysis, and this MC used the 11btunedv3.2 flux and the full 13a SuperK software.
\item We updated the value of $\sin^2\theta_{12}$ to be consistent
  with the 2013 $\nu_e$ analyses.
\end{itemize}

\subsection{Changes in the systematic effects considered by the analysis}

\begin{itemize}
\item The splines used to evaluate the uncorrelated cross-section
  systematics were regenerated using the updated SuperK MC. The
  oscillation weights used to generate this now use the nominal values
  of oscillation parameters used in the analysis, as listed in Tab.~\ref{tab:pdg2012non23}
\item The covariance matrix used to evaluate the final-state and secondary interaction
  systematics was regenerated using the updated SuperK MC.
\item The SuperK detector error matrix has been recreated using more
  SuperK atmospheric data and updated SuperK MC
  \cite{SuperkErrors2013}.
\item The SuperK energy scale systematic has changed from 2.3\% to 2.4\% \cite{SuperkEScale2013}.
\end{itemize}

\subsection{Changes in neutrino oscillation fitting}

\begin{itemize}
\item During the review process of the T2K run 1-3 result PRL \cite{T2KRun123Disapp}, it was
  discovered that the octant choice ($\theta_{23} \leq \pi/4$ or
  $\theta_{23} \geq \pi/4$) has a large effect on the contour in $\sin^2(2\theta_{23})$.
  Therefore, we now fit $\sin^2\theta_{23}$ as default
  (the ability to fit for $\sin^2(2\theta_{23})$ has been retained).
\item The hierarchy has a large effect, therefore we produce two
  results. The first using normal hierarchy; the second in inverted hierarchy.
\item In \cite{OurRun123DisappAnalysis3Flavour}, we did not consider
  the effect of uncertainties on the \emph{non-23-sector} oscillation
  parameters ($\delta_{CP}$, $\Delta m^2_{21}$, $\sin^2\theta_{13}$,
  $\sin^2\theta_{12}$). In this analysis, we consider the
  uncertainty by fitting the \emph{known} oscillation parameters ($\Delta m^2_{21}$, $\sin^2\theta_{13}$,
  $\sin^2\theta_{12}$) using the constraints listed in
  Tab.~\ref{tab:pdg2012non23}. The effect of $\delta_{CP}$ is handled
  by throwing it uniformly in toys used to created the $\Delta \chi^2$
  critical surface.
\item The 7 SuperK detector efficiency uncertainties and the 6 FSI+SI
  uncertainties have been combined into a single 6-parameter
  covariance matrix.
  This reduces the number of nuisance parameters in the fit, allowing
  the fit to run faster.
\item We now include oscillated $\nu_e \rightarrow \nu_\mu$ templates by
  reweighting the $\nu_\mu$ SuperK MC with the $\nu_e$ flux.
\end{itemize}


